%%%%%%%%%%%%%%%%%%%%%%%%%%%%%%%%%%%%%%%%
% University/School Laboratory Report
% LaTeX Template
% Version 3.1 (25/3/14)
%
% This template has been downloaded from:
% http://www.LaTeXTemplates.com
%
% Original author:
% Linux and Unix Users Group at Virginia Tech Wiki 
% (https://vtluug.org/wiki/Example_LaTeX_chem_lab_report)
%
% License:
% CC BY-NC-SA 3.0 (http://creativecommons.org/licenses/by-nc-sa/3.0/)
%
%%%%%%%%%%%%%%%%%%%%%%%%%%%%%%%%%%%%%%%%%

%----------------------------------------------------------------------------------------
%	PACKAGES AND DOCUMENT CONFIGURATIONS
%----------------------------------------------------------------------------------------

\documentclass{article}
\usepackage{float}
%\usepackage[version=3]{mhchem} % Package for chemical equation typesetting
%\usepackage{siunitx} % Provides the \SI{}{} and \si{} command for typesetting SI units
\usepackage{graphicx} % Required for the inclusion of images
\usepackage{natbib} % Required to change bibliography style to APA
\usepackage{amsmath} % Required for some math elements 
%\setlength\parindent{0pt} % Removes all indentation from paragraphs
\usepackage{vmargin}
\usepackage[dvipsnames]{xcolor}
\usepackage{amssymb}
\usepackage{multirow}
\usepackage{graphicx}
\usepackage{subfigure} % subfiguras
\usepackage{algorithm}
\usepackage[noend]{algpseudocode}
\renewcommand{\labelenumi}{\alph{enumi}.} % Make numbering in the enumerate environment by letter rather than number (e.g. section 6)
\setpapersize{A4}
\setmargins{1.75cm}       % margen izquierdo
{1.5cm}                        % margen superior
{16.5cm}                      % anchura del texto
{23.42cm}                    % altura del texto
{10pt}                           % altura de los encabezados
{1cm}                           % espacio entre el texto y los encabezados
{0pt}                             % altura del pie de página
{2cm}                           % espacio entre el texto y el pie de página
%\usepackage{times} % Uncomment to use the Times New Roman font

%----------------------------------------------------------------------------------------
%	DOCUMENT INFORMATION
%----------------------------------------------------------------------------------------

\title{Algorithm} % Title

\author{Javier Ramos, Pilar Ariza} % Author name

\date{\today} % Date for the report

\begin{document}
\begin{equation}
    \frac{\partial x_i}{\partial t} = \sum_{j \in I_H,j \neq i}\nu_i  \; {\rm e}^{ -\frac{Eb_{i\rightarrow j}}{k_BT}} \big[ x_j(1-x_i)\;{\rm e}^{\frac{\mu_j-\mu_i}{2K_BT}} - x_i(1-x_j)\;{\rm e}^{\frac{\mu_i-\mu_j}{2K_BT}} \big]
\end{equation}
\\ \\ \\
\textbf{UNITS FOR $\nu$} 
\\ \\
The second derivative of the potential with respect to the position twice, gives us the forces constant matrix of the atom \\ \\ 
 \begin{equation}\label{adp72}
 \phi_i=\frac{\partial^2 V}{\partial q_i \partial q_i} = 
     \begin{pmatrix}
\phi_{xx} & \phi_{xy} & \phi_{xz}\\
\phi_{yx} & \phi_{yy} & \phi_{yz}\\
\phi_{zx} & \phi_{zy} & \phi_{zz}
\end{pmatrix}\;\;[\frac{eV}{\r{A}^2}]
 \end{equation}
 \\
using the analogy of a spring, we can consider
\\
\begin{equation}\label{adp73}
    k=\sum_i \phi_{ii}, \; \;\; \;\; \;\; \;  \nu_i=\sqrt{\frac{k}{3m_i}},
\end{equation}

In order to work with frequency unit we have to make a unit conversion. If we recall the Eq. (\ref{adp73}), the initial units are
\\
\begin{equation} \label{adp80}
    \sqrt{\frac{eV/\r{A}^2}{u}},
\end{equation}
\\
the first unit conversion we do is $1 eV/\r{A}^2 =$ $16.025 J/m^2 $, $1J = 1 N \cdot m$ and $1 u = $ $1.66^{-24}g$. Putting this together we have
\\
\begin{equation}\label{adp81}
    \sqrt{\frac{16.025 N \cdot m/{m}^2}{1.66^{-24}g}}.
\end{equation}
\\
The second unit conversion we do is $1N = 10^3 g \cdot m /s^2$ and we finally obtain
\\
\begin{equation}\label{adp82}
    \sqrt{\frac{(eV\cdot \frac{1.60218 \cdot 10^{-19}J}{1 eV}) /(\r{A}^2 \cdot \frac{10^{-20}m^2}{1 \r{A}^2})}{u \cdot \frac{1.660538921 \cdot 10^{-27} kg}{1 u}}} = 98.25 \cdot 10^{12} \frac{1}{s} 
\end{equation}
\\ \\
\textbf{UNITS FOR $\mu$} 
\\ \\
\begin{equation}\label{mu}
    \mu = \gamma * K_BT
\end{equation}

\begin{equation} \label{gamma}
    \gamma = \frac{3}{2} + log\frac{x}{1-x}+\frac{1}{k_bT}[\frac{1}{eV}]\frac{\partial V}{\partial x}[eV]
\end{equation}
In Eq. (\ref{gamma}) we see that $\gamma$ is dimensionless, from Eq. (\ref{mu}) we know that the dimensions for $\mu$ are eV.
\\ \\
\textbf{UNITS FOR $f_i$} 
\\ \\
\begin{equation}
    f_i = \frac{\partial V}{\partial q_i} \;\; [\frac{eV}{\r{A}}]
\end{equation}
\\ \\
\textbf{UNITS FOR $V$} 
\\ \\
\begin{equation}
V [eV]
\end{equation}
\\ \\
\textbf{VALUES FOR $\nu_i$} 
\\ \\
\begin{equation}
     k=\sum_i \phi_{ii}=(1.189,...,6.168) \;\; [\frac{eV}{\r{A}^2}]
\end{equation}
\\
\begin{equation}
    \frac{k}{3m}=(0.389,...,2.039)\;\; [\frac{eV}{\r{A}^2u}]
\end{equation}
\\
\begin{equation}
    \nu_i=\sqrt{\frac{k}{3m_i}}=(0.623,...,1.47) [\sqrt{\frac{eV}{\r{A}^2u}}] = (61.21,...,144.42)  \cdot 10^{12} [\frac{1}{s}]
\end{equation}
\\ \\ \\
$\nu_i = 100*10^{12} s^{-1}.$
\\
${E^b}_{i\rightarrow j}=0.1 eV$
\\
$K_B = 8.617333262 \cdot 10^{-5} eV/K$
\\
$T = 350K$
\\
$x_i = 0.7$
\\
$x_j = 0.3$
\\
$\mu_i = -3 eV$
\\
$\mu_j = -3.7 eV$
\\ \\ \\
\begin{equation}\nonumber
k_BT = 0.03016 eV
\end{equation}

\begin{equation}\nonumber
    -\frac{{E^b}_{i\rightarrow j}}{K_BT}=-3.31565
\end{equation}

\begin{equation}\nonumber
    {\rm e}^{-\frac{{E^b}_{i\rightarrow j}}{K_BT}}=0.03631
\end{equation}

\begin{equation}\nonumber
    \nu_i  \;{\rm e}^{-\frac{{E^b}_{i\rightarrow j}}{k_BT}}=4.35*10^{12} \; s^{-1}
\end{equation}

\begin{equation}\nonumber
    x_j(1-x_i)=0.3·(1-0.7)=0.09
\end{equation}

\begin{equation} \nonumber
   {\rm e}^{\frac{\mu_j-\mu_i}{2K_BT}} = {\rm e}^{\frac{-3.7 \; eV- (- 3) \; eV}{2*0.03016eV}}={\rm e}^{\frac{-0.7 \; eV}{0.06032 \; eV}}=9.12 \cdot 10^{-6}
\end{equation}

\begin{equation} \nonumber
   x_j(1-x_i){\rm e}^{\frac{\mu_j-\mu_i}{2K_BT}} =0.09 \cdot 8.57 \cdot 10^{-6} = 8.21 \cdot 10^{-7}
\end{equation}

\begin{equation}\nonumber
    x_i(1-x_j)=0.7 \cdot (1-0.3)=0.49
\end{equation}

\begin{equation} \nonumber
   {\rm e}^{\frac{\mu_i-\mu_j}{2K_BT}} = {\rm e}^{\frac{-3eV- (- 3.7) \; eV}{2*0.03016 \; eV}}={\rm e}^{\frac{0.7 \; eV}{0.06032 \; eV}}=109619
\end{equation}

\begin{equation} \nonumber
   x_i(1-x_j) {\rm e}^{\frac{\mu_i-\mu_j}{2K_BT}} =0.49 \cdot 109619=53713.31
\end{equation}

\begin{equation}\nonumber
    \big[ x_j(1-x_i)\;{\rm e}^{\frac{\mu_j-\mu_i}{2K_BT}} - x_i(1-x_j)\;{\rm e}^{\frac{\mu_i-\mu_j}{2K_BT}} \big] = 8.21 \cdot 10^{-7} - 53713.31 = - 53713.31
\end{equation}

\begin{equation}
    \frac{\partial x_i}{\partial t} = \sum_{j \in I_H,j \neq i}\nu_i  {\rm e}^{ -\frac{{E^b}_{i\rightarrow j}}{K_B T}} \big[ x_j(1-x_i)\;{\rm e}^{\frac{\mu_j-\mu_i}{2K_B T}} - x_i(1-x_j)\;{\rm e}^{\frac{\mu_i-\mu_j}{2K_B T}} \big]= 4.35 ^10^{12}\; s^{-1} \cdot (- 53713.31)=233652*10^{12} \; s^{-1}
\end{equation}
\\ \\ \\
\textbf{NEW MASTER EQUATION}
\\ \\
\begin{equation}
    \frac{\partial x_i}{\partial t}= \sum_{j \in I_H,j \neq i} \bigg( \nu_j{\rm e}^{-\beta {E}_{j\rightarrow i}} x_j(1-x_i)  - \nu_i{\rm e}^{-\beta {E}_{i\rightarrow j}} x_i(1-x_j) \bigg)
\end{equation}
\\ \\
$\nu_i = 120 *10^{12}s^{-1}.$
\\
$\nu_j = 100*10^{12} s^{-1}.$
\\
${E}_{i\rightarrow j}=0.1 eV$
\\
${E}_{j\rightarrow i}=0.3 eV$
\\
$K_B = 8.617333262 \cdot 10^{-5} eV/K$
\\
$T = 350K$
\\
$x_i = 0.7$
\\
$x_j = 0.3$
\\ \\
\begin{equation}
    \beta = \frac{1}{K_BT}=33.157 eV
\end{equation}
\\
\begin{equation}
\beta * {E}_{j\rightarrow i}=9.947
\end{equation}
\\
\begin{equation}
    {\rm e}^{-\beta {E}_{j\rightarrow i}} = {\rm e}^{-9.947}=4.78*10^{-5}
\end{equation}
\\
\begin{equation}
    \nu_j{\rm e}^{-\beta {E}_{j\rightarrow i}} x_j(1-x_i) =  100*10^{12}[s^{-1}]*4.78*10^{-5} * 0.3(1-0.7)=4.30*10^{8}[s^{-1}]
\end{equation}
\\
\begin{equation}
\beta * {E}_{i\rightarrow j}=3.3157
\end{equation}
\\
\begin{equation}
     {\rm e}^{-\beta {E}_{i\rightarrow j}} = {\rm e}^{-3.3157}=0.0363
\end{equation}
\\
\begin{equation}
    \nu_i{\rm e}^{-\beta {E}_{i\rightarrow j}} x_i(1-x_j) = 120*10^{12}[s^{-1}]*0.0363 * 0.7(1-0.3)=2.13*10^{12}[s^{-1}]
\end{equation}
\\
\begin{equation}
    \bigg( \nu_j{\rm e}^{-\beta {E}_{j\rightarrow i}} x_j(1-x_i)  - \nu_i{\rm e}^{-\beta {E}_{i\rightarrow j}} x_i(1-x_j) \bigg)=(4.30*10^{8}-2.13*10^{12})[s^{-1}]=-2.13*10^{12}(\frac{1}{s})
\end{equation}
\end{document}